% ============================================================
%  APPENDIX — EXPERIMENT 3 FULL DETAILS  (training & transfer)
%  Place after the Experiment 2 appendix.
% ============================================================

\section{Experiment 3: Controlled Relationship Transfer}
\label{app:exp3}

As the controlled transfer step, Experiment~3 asks whether outcome-based training
learns to select and apply an episode-matched relationship beyond its training
setting. The primary Qwen3-8B endpoint jointly changes the domain and response
mechanism; Qwen3-4B and Llama-3.1-8B provide scale and architecture comparisons.

\subsection{Design and Evaluation}
\label{app:exp3-coin-structural}

\subsubsection{Data-generating processes}

Training uses only Structure~A, the direct relationship from Experiment~2. On
the displayed scale of domain $d$,
\begin{equation}
  y_{t+1}=b_d+g_jx_t+\epsilon_t,\qquad
  \epsilon_t\sim\mathcal N(0,\sigma_d^2).
  \label{eq:exp3-coin-a}
\end{equation}
A pulse acts in the next period and has no residual effect under zero subsequent
input. The unseen evaluation mechanism, Structure~B, is
\begin{align}
  m_{t+1}&=\rho m_t+g_jx_t,\\
  y_{t+1}-b_d&=\phi(y_t-b_d)+\lambda m_{t+1}+\epsilon_t.
  \label{eq:exp3-coin-b}
\end{align}
This is a linear-Gaussian latent-state process in the sense of
\citet{kalman1960new}: it introduces a persistent mediator and outcome carry-over,
so a single pulse has both immediate and delayed effects. Model prompts disclose neither equation,
parameter values, internal variable, nor structure label.

Coin City uses $b=50$, range $[0,100]$, driver pool
$\{-12,-8,-4,0,4,8,12\}$, and $\sigma=1.6$. Held-out Coin Harbor uses
$b=180$, range $[40,320]$, driver pool
$\{-24,-16,-8,0,8,16,24\}$, and $\sigma=4.5$, together with new terminology,
units, and qualitative descriptions. Evaluation gains follow
$g\sim\mathrm{Uniform}(.22,.92)$. Structure~B independently samples
$\rho\sim\mathrm{Uniform}(.48,.74)$,
$\phi\sim\mathrm{Uniform}(.18,.42)$, and
$\lambda\sim\mathrm{Uniform}(.68,.94)$. A reference and its matched target share
these parameters but have independent inputs and noise.

\begin{center}
\begin{minipage}{\linewidth}
\centering
\captionof{table}{\textbf{Exact Coin City--Coin Harbor domain contrast.} Both
surfaces instantiate the same Structure~A and B simulators and task template; Coin
Harbor changes the semantic setting, time unit, numeric scale, interventions, noise,
and pathway descriptions.}
\label{tab:exp3-domain-contrast}
\footnotesize
\setlength{\tabcolsep}{3pt}
\renewcommand{\arraystretch}{1.05}
\begin{tabular}{p{.20\linewidth}p{.36\linewidth}p{.36\linewidth}}
\toprule
\textbf{Field} & \textbf{Coin City (train and test)} & \textbf{Coin Harbor (test only)} \\
\midrule
Driver $\to$ outcome & Net campaign news $\to$ candidate support & Net shipping
bulletin $\to$ cargo-flow index \\
Time unit & Week & Tide cycle \\
Resting level / support & 50 / $[0,100]$ & 180 / $[40,320]$ \\
Observed-input pool & $\{-12,-8,-4,0,4,8,12\}$ &
$\{-24,-16,-8,0,8,16,24\}$ \\
Evaluation shocks & $\{-12,-6,0,6,12\}$ & $\{-24,-12,0,12,24\}$ \\
Observation-noise SD & 1.6 support points & 4.5 cargo-index points \\
Structure A description & Synchronized citywide alerts; reactions occur mostly in
the same week & Synchronized fleetwide radio; responses occur mostly in the same
tide cycle \\
Structure B description & Neighborhood conversations and local organizations;
reactions build gradually and persist & Dockside crews and harbor groups; responses
build gradually and persist \\
\bottomrule
\end{tabular}
\end{minipage}
\end{center}

\subsubsection{Training episodes and matched control}

Each training prompt contains two 14-case Coin City reference systems, both
governed by Equation~\ref{eq:exp3-coin-a}. One draws
$g\in[.18,.46]$ and the other $g\in[.68,.98]$; their order is randomized.
The target shares exactly one reference gain, selected in balanced fashion, and
shows $k\in\{0,2,4,8\}$ independently generated cases. Its national/local
background description identifies the matching weak or strong response as in
Experiment~2. There are 1,200 rows at each $k$, or 4,800 per training run.

The episode-matched and population-prior datasets have byte-identical prompts.
Episode-matched reward uses the target's noise-free forecasts. For the control, targets are
deranged within each $k$ stratum with a fixed permutation: no row keeps its own
target, while the complete multiset of target vectors is exactly preserved.
Consequently the control can learn the output format and population distribution
but cannot improve reward by using that row's references or target history.
The Qwen3-4B structureless diagnostic additionally scrambles displayed drivers
and rewards a flat intervention response. It isolates task-format learning and is
excluded from the episode-matched-versus-prior estimand.

% Worked example for the Experiment 3 training policies.
%
% Panel A is the real content of training episode train:1:k2, regenerated from
% exp3_training_transfer/coin_city_structural/make_dataset.py (_training_base(1));
% backgrounds are abridged from the rendered prompt.
%
% Panel B draws the prompt-to-key pairing that defines the three trained arms.
% The population-prior control is a cyclic shift within an evidence-depth
% stratum (np.roll under seed 0xC01AC17 in build_training), so it is drawn as a
% shift with one wrap-around edge: every key is still used exactly once, but no
% prompt keeps its own.
\begin{center}
\begin{minipage}{\linewidth}
\captionof{table}{\textbf{One training episode, and what each policy is rewarded against.}
Panel~A is the prompt for episode \texttt{train:1:k2}, abridged from the rendered
text. Panel~B is the prompt--key pairing used during training, drawn for four episodes
at one evidence depth: $P_i$ is episode $i$'s prompt and $\mathbf y^\star_i$ its
ten-answer key. Episode-matched and population-prior training use byte-identical prompts
and the same multiset of keys; only the pairing differs. The control shifts
cyclically across all 1,200 episodes at a depth, so its key is a real key drawn
from the same distribution but never this prompt's, and no reward can be earned by
reading this row's references. Base is untrained and receives no supervision.}
\label{tab:exp3-training-policies}
\scriptsize
\setlength{\tabcolsep}{2.5pt}
\renewcommand{\arraystretch}{1.12}
% The two paper builds have different measures; shrink a panel only if its
% natural width would overrun the text block.
\providecommand{\fitwidth}[1]{%
  \resizebox{\ifdim\width>\linewidth\linewidth\else\width\fi}{!}{#1}}

% ---- Panel A: the shared prompt ------------------------------------------
\noindent\colorbox{black!4}{%
\begin{minipage}{\dimexpr\linewidth-2\fboxsep\relax}%
\fitwidth{%
\begin{tabular}{@{\ }l@{\ \ }l@{\ }}
\multicolumn{2}{@{\ }l@{\ }}{\textbf{A. Shared prompt.} Coin City; resting level
$50.0$, bounded to $[0,100]$; no equation or coefficient is shown.} \\[1.5pt]
\textsc{reference 1} & ``\dots immediate national alerts \dots react
\textbf{strongly}''\quad 14 wk: $-8.0\!\to\!43.8$, $-12.0\!\to\!40.9$, \dots,
$0.0\!\to\!51.6$ \\
\textsc{reference 2} & ``\dots local civic briefings \dots react
\textbf{weakly}''\quad 14 wk: $+8.0\!\to\!52.6$, $+4.0\!\to\!51.3$, \dots,
$-12.0\!\to\!46.9$ \\
\textsc{target} ($k{=}2$) & ``\dots immediate national alerts \dots react
\textbf{strongly}''\quad $-8.0\!\to\!42.3$, $+4.0\!\to\!53.9$ \\
\textsc{query} & ten counterfactuals A--J: shock $\in\{-12,-6,0,6,12\}$ $\times$
horizon $\in\{1,3\}$; return \texttt{\{"forecasts":[A,\dots,J]\}} \\
\end{tabular}}%
\end{minipage}}

\vspace{6pt}
\noindent\textbf{B. How prompts are paired with answer keys during training.}
\vspace{2pt}

% ---- Panel B: the prompt-to-key wiring ------------------------------------
\definecolor{wireInk}{HTML}{3F4550}
\definecolor{wireMute}{HTML}{8A93A3}
\definecolor{wireKeyFill}{HTML}{E8F2FB}
\definecolor{wireKeyEdge}{HTML}{4F7EA8}
\noindent\fitwidth{%
\begin{tikzpicture}[
  x=1cm, y=1cm,
  pchip/.style={rounded corners=2pt, minimum width=0.95cm, minimum height=0.26cm,
    inner sep=0pt, draw=wireMute!60, fill=black!4, line width=.5pt,
    font=\scriptsize, text=wireInk},
  schip/.style={pchip, dash pattern=on 1.3pt off 1.1pt},
  kchip/.style={rounded corners=2pt, minimum width=0.95cm, minimum height=0.26cm,
    inner sep=0pt, draw=wireKeyEdge, fill=wireKeyFill, line width=.5pt,
    font=\scriptsize, text=wireInk},
  flatbox/.style={kchip, minimum width=1.60cm, minimum height=0.30cm,
    draw=wireMute!70, fill=black!4},
  wire/.style={-{Stealth[length=3pt,width=2.4pt]}, draw=wireInk!65,
    line width=.5pt, shorten >=.5pt, shorten <=.5pt},
  ptitle/.style={anchor=south, font=\scriptsize\bfseries, text=wireInk},
  pnote/.style={anchor=north, font=\tiny\itshape, text=wireMute},
]
\node[ptitle] at (1.55,0.12) {Episode-matched};
\node[ptitle] at (6.65,0.12) {Population prior};
\node[ptitle] at (11.72,0.12) {Structureless};

% --- Episode-matched: each prompt keeps its own key --------------------------------
\foreach \r/\y in {1/-0.30, 2/-0.66, 3/-1.02, 4/-1.38}{
  \node[pchip] (cp\r) at (0.50,\y) {$P_\r$};
  \node[kchip] (ck\r) at (2.60,\y) {$\mathbf y^\star_\r$};
  \draw[wire] (cp\r.east) -- (ck\r.west);}
\node[pnote] at (1.55,-1.85) {its own key};

% --- Population prior: cyclic shift within an evidence depth ---------------
% The wrap edge is routed under the chips and up outside the key column so it
% cannot clip a chip it does not connect to.
\foreach \r/\y in {1/-0.30, 2/-0.66, 3/-1.02, 4/-1.38}{
  \node[pchip] (pp\r) at (5.60,\y) {$P_\r$};
  \node[kchip] (pk\r) at (7.70,\y) {$\mathbf y^\star_\r$};}
\foreach \a/\b in {1/2, 2/3, 3/4}{
  \draw[wire] (pp\a.east) -- (pk\b.west);}
\draw[wire, rounded corners=3.5pt]
  (pp4.south) -- (5.60,-1.68) -- (8.50,-1.68) -- (8.50,-0.30) -- (pk1.east);
\node[pnote] at (6.65,-1.85) {every key used once, none its own};

% --- Structureless: scrambled prompt, one flat key -------------------------
\foreach \r/\y in {1/-0.30, 2/-0.66, 3/-1.02, 4/-1.38}{
  \node[schip] (sp\r) at (10.70,\y) {$\tilde P_\r$};}
\node[flatbox] (sflat) at (12.75,-0.84) {$(50,\dots,50)$};
\foreach \r in {1,2,3,4}{\draw[wire] (sp\r.east) -- (sflat.west);}
\node[pnote] at (11.72,-1.85) {flat key; displayed drivers permuted};
\end{tikzpicture}}
\end{minipage}
\end{center}


\subsubsection{Evaluation interface and qualitative hint}

Every evaluation prompt contains a complete 14-period Structure~A reference, a
complete 14-period Structure~B reference, and a target governed by one of them.
For Coin City the two reference descriptions are:
\begin{quote}\small
Residents receive campaign developments through synchronized citywide alerts.
Reactions occur mostly during the same week as the news.

Campaign developments usually spread through neighborhood conversations and
local organizations. Reactions may build gradually and persist after the
original news has faded.
\end{quote}
Coin Harbor uses the matched fleetwide-radio versus dockside-network descriptions.
The target receives the correct description, no description, or the other
pathway's description. The three prompt variants reuse the exact reference and
target trajectories, scenario grid, and gold forecasts; the manifest records a
common numeric hash for every cue triplet.

Each task requests ten counterfactuals: shocks
$\{-2u_d,-u_d,0,u_d,2u_d\}$ at horizons one and three, with $u_d=6$ in Coin City
and $u_d=12$ in Coin Harbor. The shock occurs in the next period and all later
inputs are zero. The response vector consists of the eight nonzero-shock minus
zero-shock contrasts. The full endpoint contains
\[
2\ {\rm domains}\times2\ {\rm structures}\times3\ {\rm hints}\times
4\ k{\rm\ values}\times30\ {\rm episodes}=1,440
\]
paired prompts. Training and evaluation seed bands begin at 81,000,000 and
91,000,000, respectively.

\paragraph{Exact task template and prompt variants.}
\label{app:exp3-coin-prompts}
The following is the task string emitted by the prompt generator before the
model-specific chat wrapper is applied. Bracketed uppercase fields below denote
row-specific substitutions; they are not shown to the model. Line wrapping is
typographic, while the blank-line-separated sections and all fixed wording are
part of the prompt.

\begin{quote}
\begingroup\ttfamily\scriptsize\raggedright
Forecasting task in [DOMAIN].

\smallskip
The resting level of [OUTCOME] is [BASELINE]. Values are bounded to
[LOWER, UPPER].

\smallskip
The reference systems are complete examples. The target system may share a
stable response pattern with one reference. Infer from the information shown;
no response equation or coefficient is provided.

\smallskip
REFERENCE SYSTEM 1

Background: [REFERENCE 1 CONTEXT]

[REFERENCE 1 TABLE: period, driver, and 14 observed outcomes]

\smallskip
REFERENCE SYSTEM 2

Background: [REFERENCE 2 CONTEXT]

[REFERENCE 2 TABLE: period, driver, and 14 observed outcomes]

\smallskip
TARGET SYSTEM

[TARGET BACKGROUND, WHEN PRESENT]

[TARGET TABLE WITH $k$ OBSERVATIONS, OR THE EXACT NO-OBSERVATION SENTENCE]

\smallskip
COUNTERFACTUAL FORECASTS

For each scenario, apply the stated [DRIVER] in the next [PERIOD] and then set it
to 0 in later periods. Forecast [OUTCOME] at the stated horizon.

\smallskip
A: shock=[VALUE]; horizon=1\quad $\ldots$\quad
E: shock=[VALUE]; horizon=1

F: shock=[VALUE]; horizon=3\quad $\ldots$\quad
J: shock=[VALUE]; horizon=3

\smallskip
Return only strict JSON with exactly ten finite numbers in scenario order:

\{"forecasts":[A,B,C,D,E,F,G,H,I,J]\}
\endgroup
\end{quote}

The renderer realizes the variants as follows. Training always includes the
correct strong/weak background and uses the same input string for episode-matched and
population-prior training; only the hidden reward target differs. At evaluation,
\texttt{correct} inserts the target pathway's description, \texttt{none} omits
the background line entirely, and \texttt{misleading} inserts the other
pathway's description. Every cue triplet otherwise has identical tables,
interventions, and targets. If $k=0$, the target block contains the exact
sentence \emph{No target-system observations are available yet.}; otherwise it
contains the same table format as the references with the first $k$ rows. The
Qwen3-4B \texttt{structureless} diagnostic retains this shell but swaps every
displayed $-12.0\leftrightarrow12.0$ and $-8.0\leftrightarrow8.0$ token through
a collision-safe substitution. Untrained bases receive the same held-out
prompts as trained models. Thus cue, reward-arm, structureless, and base
comparisons do not introduce an unreported instruction change.

\subsubsection{Estimands, roster, and validation criteria}

The reported factorial crosses Qwen3-4B, Qwen3-8B, and Llama-3.1-8B, two reward
arms, and three training seeds, for 18 runs. Three Qwen3-4B structureless runs
and three untrained endpoints provide diagnostic comparisons.
The primary estimand is population-prior minus episode-matched response MAE on
Structure~B in Coin Harbor with the correct hint, averaging the four $k$ values
equally. Positive values favor episode-matched training. The same estimator defines
factorial contrasts for Structure~A/Coin Harbor (domain transfer),
Structure~B/Coin City (structural transfer), and the in-distribution cell.
Paired absent-minus-correct and misleading-minus-correct contrasts measure hint
use. Their paired $k=8$ minus $k=0$ difference tests whether observations override
a misleading prior; negative values indicate a shrinking misleading-hint penalty.

Every reported round-300 endpoint uses five draws per prompt at temperature $.7$;
temperature-zero checkpoints serve only as online training monitors.

Each episode's score is the mean of its five per-draw response MAEs, so ``averaging
draws'' averages errors and not forecasts. The hierarchical bootstrap
resamples the three training seeds at the top level and paired episodes within
each selected seed. Parse failures receive the full observable-range error
penalty rather than being discarded. That penalty is the clip width---280 points
in Coin Harbor, roughly fifty times a typical error---so it carries large
leverage: a single unparsed draw of 150 raises the Qwen3-8B base joint-cell MAE
from $5.26$ to $5.72$, which is most of that arm's reported gap to matched
training. Both trained arms parse completely in that cell, so the primary
estimand is unaffected; the base comparison is not.
The launch ledger contains all 24 endpoints and 172,800 stochastic rows;
172,732 parse (99.96\%), and every endpoint has at least 99.82\% coverage.

\paragraph{Sensitivity of the intervals to three training seeds.}
\label{app:exp3-coin-cluster-sensitivity}
The hierarchical interval resamples three top-level clusters, and a nonparametric
cluster bootstrap at $G=3$ is known to understate between-cluster uncertainty
\citep{cameron2008bootstrap}. We therefore report a deliberately conservative
alternative alongside it: a Student-$t$ interval on the three seed-level means with
two degrees of freedom, which uses the same paired episode differences but charges
the full small-sample penalty for estimating a variance from three numbers. Under
that alternative none of the twelve transfer-cell intervals excludes
zero, the primary Qwen3-8B contrast becoming $.251$ $[-.259,.762]$ and the Qwen3-4B
comparator $.319$ $[-.120,.758]$. Nine of the twelve intervals that exclude zero
under the hierarchical bootstrap therefore do not survive the $t$ alternative, and we
do not claim that any single interval in this experiment is robust to the choice of
small-cluster procedure.

What does survive is the sign pattern, which requires no variance estimate. Across
the four transfer cells and three training seeds there are twelve seed-level
contrasts per model. All twelve favor episode-matched training for Qwen3-8B and all
twelve for Qwen3-4B, so all 24 Qwen seed-by-cell contrasts point the same way;
these reuse training seeds across cells and are not independent replications. The corresponding count for
Llama-3.1-8B is three of twelve. We take the two-Qwen sign consistency, rather than
any interval, as the principal evidence for the episode-matching effect, and we
read the interval estimates as descriptive of magnitude.

Before training, the dataset passed 17 fail-closed checks: exact row counts and factorial
balance; byte-identical episode-matched/population-prior prompts and held-out datasets; exact
within-$k$ target marginals with no fixed reward links; absence of Structure~B
from training prompts and metadata; cue-triplet numeric identity; train/evaluation
identifier disjointness; hidden-equation leak scans; strict gold round trips; and
oracle separation from the blind pathway mixture. Mean A--B response separation
is 2.65 points in Coin City and 4.80 in Coin Harbor. Across both tokenizers,
the longest prompt is 1,071 tokens against the configured 2,304-token limit.

We estimate population-prior-minus-episode-matched intervals in all four transfer
cells, paired intervals for both cue penalties at every $k$, and the paired
$k=8$ minus $k=0$ change in misleading-minus-correct error to evidence override.
The protocol archive records versioned input, analysis, and output hashes. Full
optimization and compute parameters are in
Section~\ref{app:exp3-complete-training-setup}.

\subsubsection{Stochastic scale and architecture comparators}
\label{app:exp3-coin-qwen-results}

\paragraph{Qwen3-4B scale replication.}
\label{app:exp3-coin-qwen4b}
The Qwen3-4B comparison includes three seeds in each of the episode-matched,
population-prior, and structureless arms. Each is evaluated on the same 1,440
tasks at round 300 with five stochastic draws per task; the untrained base uses
the same endpoint. Parse rates are 100\% for every trained arm
and 99.97\% for the stochastic base.

Joint-cell MAE is $6.15$, $6.76$, $4.81$, and $4.49$ for base,
structureless, population-prior, and episode-matched training, respectively.
The prior-minus-episode-matched difference is $.319$ $[.134,.489]$;
episode-matched training has the lowest point estimate in all four cells, and
all four intervals exclude zero. The Qwen3-8B main estimate is likewise positive,
$.251$ $[.058,.470]$, giving the same interval-excluding-zero result at both
tested Qwen sizes.

Averaged over evidence depths, omitting the hint raises episode-matched response
MAE by $.151$ $[.079,.235]$, and a misleading hint raises it by $.313$
$[.228,.415]$. The misleading-hint penalty is smaller at $k=8$ than at $k=0$,
but the evidence-override contrast $-.116$ $[-.309,.100]$ includes zero.

\paragraph{Llama-3.1-8B architecture comparator.}
In the joint-transfer cell the prior-minus-episode-matched training contrast is
positive, $.420$ $[-.115,.958]$, but its interval spans zero. The three seed-level
effects are $+.762$, $+.622$, and $-.125$, so this result does not establish a
robust cross-architecture training advantage.

The joint cell is also the most favorable of the four Llama cells, and
reporting it alone would misrepresent the factorial, so we give all four. The
contrasts are $-.351$ $[-.680,-.130]$ in the in-distribution cell, $-.167$
$[-.862,.306]$ under structure-only shift, $-.581$ $[-.868,-.243]$ under
domain-only shift, and $+.420$ $[-.115,.958]$ under joint shift. In the two
Structure~A cells the population-prior arm is therefore better than the
episode-matched arm by an interval that excludes zero, and only three of the twelve
Llama seed $\times$ cell contrasts favor episode-matched training. On this
architecture the factorial runs against the hypothesis in the cells
whose response mechanism was seen during training, and recovers a positive point
estimate only where both shifts apply. We report this as a failure of the
episode-matching effect to replicate on Llama-3.1-8B rather than as a
directionally-consistent-but-underpowered result, and the manuscript's transfer
claim rests on the two Qwen models.

Llama nevertheless uses the qualitative
pathway cue at the stochastic endpoint:
absent-minus-correct MAE is $.219$ $[.123,.313]$ and
misleading-minus-correct is $.210$ $[.092,.320]$. Its misleading-hint penalty
shrinks from $k=0$ to $k=8$ by $-.621$ $[-.957,-.323]$. All three cue-related
intervals exclude zero. These contrasts indicate sensitivity to the supplied
mechanism cue and reduced misleading-cue effects as observations accumulate.

\begin{center}
\begin{minipage}{\linewidth}
\centering
\captionof{table}{\textbf{Stochastic Coin City appendix comparators in the
joint-transfer cell.} The Qwen3-8B primary result remains in the main body.
Entries are response MAE under the correct hint after averaging five draws per
task; lower is better. Positive prior-minus-episode-matched contrasts favor
episode-matched training. Intervals use the paired seed-first
bootstrap with a fixed cell-specific bootstrap seed.}
\label{tab:exp3-coin-complete-roster}
\scriptsize
\setlength{\tabcolsep}{3pt}
\resizebox{\linewidth}{!}{% generated by coin_city_structural/make_paper_outputs.py -- do not edit
\begin{tabular}{lrrrr}
\toprule
Model & Base & Prior & Episode-matched & $\Delta$ [95\% CI] \\
\midrule
Qwen3-4B & 6.15 & 4.81 & 4.49 & 0.32 [0.13,0.49] \\
Llama-3.1-8B & 9.43 & 4.69 & 4.27 & 0.42 [-0.12,0.96] \\
\bottomrule
\end{tabular}
}
\end{minipage}
\end{center}

\section{Experiment 4: Historical-Market Training and Held-Out Evaluation}
\label{app:exp4}
\label{app:exp3b-protocol}

Experiment~4 asks whether base-relative gains extend from the controlled setting
to events and normalized question families unseen during training. It uses
proper-score training, a point-in-time cutoff, and family-disjoint held-out data.

\subsection{Sample Construction and Leakage Controls}

\paragraph{Source archive and sample.}
The source export\footnote{%
  \micon{huggingface.png}\;%
  \href{https://huggingface.co/datasets/anonymous-review/polymarket-full-market-data}%
  {Full Polymarket archive}. Link anonymized for review.%
} has 1,788,703 rows and SHA-256
\texttt{5697f767\ldots f638e7d} (the full digest is stored in the manifest).
The builder scans the entire file and fails on checksum drift. The calendar
splits contain 1,736 train, 512 development, and 1,024 test markets, with YES rates
.311, .275, and .301. Train contains 991 distinct events and 1,456 normalized
question templates; development contains 292 and 453; test contains 493 and 866.
The selected archive rows have no official series identifiers, so event and
template components supply the operative family groups.

\paragraph{Eligibility and point-in-time censoring.}
We admit literal Yes/No markets resolved by an extreme final token price. A
boundary-aware domain rule selects political, governmental, macroeconomic, and
geopolitical questions while explicitly rejecting culture, awards, sports,
esports, and crypto categories. This rule is fixed before model training. One task
is built 30 days before the earliest archived terminal timestamp. At least four
distinct UTC dates must have genuine pre-cutoff YES candles, the latest candle must
be no more than seven days stale, and its price must remain in $[.02,.98]$. We keep
at most the latest 14 dates. Empty, carry-forward, and post-cutoff candles are
rejected. The prompt omits the outcome, final prices and volume, market status,
realized terminal time, and the tags used for domain selection.

The archive records timestamped price history but not revisions to static question
or resolution-rules text. We therefore treat those strings as time-invariant. This
is a residual point-in-time limitation: the audit verifies that no timestamped or
outcome field crosses the cutoff, but cannot determine whether an archived rule
changed between the forecasting cutoff and resolution.

\paragraph{Family-disjoint temporal split.}
Before applying sample caps, every development row sharing an event, available
series, or normalized template with train is removed; the same rule removes test
rows matching train or development. Numeric values and month names are normalized
in template keys. The final intersection is zero for each pair of splits. Sampling
within a split is a stable hash of task identifier and protocol version, so row
order in the export cannot change membership.

\subsection{Training and Held-Out Evaluation}

\paragraph{Matched update budget.}
Training comprises 1,736 unique markets arranged in a deterministic, balanced
schedule of 4,800 presentations, matching Experiment~3's 300-round update budget.
Every market appears twice and 1,328 appear a third time. Schedule order and task
files are checksum-verified. With batch size 16, this gives exactly 300
rollout-and-update rounds (2,400 AdamW steps). The manifests and tables distinguish
unique markets from presentations. The longest Qwen3-4B wrapped prompt contains
1,012 tokens against the 1,792-token training limit. Data, schedule, reward,
optimizer, update count, seeds, and scoring rule are held constant across models;
each uses its native chat template.
Shared optimizer, batching, and runtime details appear in
Section~\ref{app:exp3-complete-training-setup}.

\begin{center}
\begin{minipage}{\linewidth}
\centering
\captionof{table}{Experiment~4 training configuration. Held-out data are absent
from training and evaluated once per model after all three training seeds.}
\label{tab:exp3b-hparams}
\scriptsize
\begin{tabular}{p{.34\linewidth}p{.59\linewidth}}
\toprule
\textbf{Quantity} & \textbf{Value} \\
\midrule
Model roster / adapter & Qwen3-4B-Instruct-2507, Qwen3-8B, and
Llama-3.1-8B-Instruct; LoRA rank 32, $\alpha=64$ \\
Algorithm & Dr.~GRPO; $\beta=0$; AdamW, lr $10^{-6}$ \\
Unique train / presentations & 1,736 / 4,800 balanced repetitions \\
Rounds / seeds & 300; $42,43,44$ \\
Rollout & 16 prompts, 8 samples each; temperature 1.3 \\
Output & one JSON probability; 128 new-token limit \\
Reward & $1-(\hat p-y)^2$ only \\
Chat formatting & Qwen3-4B uses its frozen ChatML wrapper; Qwen3-8B uses
its non-thinking template; Llama uses its native instruct template. \\
Development & 512 tasks; temperature-zero evaluation every 25 rounds \\
Held-out evaluation & 1,024 tasks; temperature zero; one evaluation per model \\
\bottomrule
\end{tabular}
\end{minipage}
\end{center}

\subsection{Held-Out Results}
\label{app:exp3b-results}
\paragraph{Decision rule.}
The primary learned contrast is the mean loss over three trained seeds minus base
loss. Invalid JSON receives Brier loss one. Intervals use 5,000 bootstrap
resamples of connected event/template components.

All Qwen3-8B test forecasts parse. Its three-seed trained mean obtains Brier
.11642 versus .12867 for the base, a difference of $-.01225$ with 95\% interval
$[-.01948,-.00602]$. Every training seed improves on the base. The interval
excludes zero, supporting transfer to held-out events for Qwen3-8B.

\begin{center}
\begin{minipage}{.94\linewidth}
\centering
\captionof{table}{\textbf{Qwen3-8B on the 1,024-market held-out set.}
The mean-row interval resamples connected event/template families; lower is
better.}
\label{tab:exp4-endpoint-results}
\scriptsize
\setlength{\tabcolsep}{3.5pt}
\resizebox{\linewidth}{!}{% generated by polymarket/scripts/render_model_extension_table.py -- do not edit
\begin{tabular}{lrrr}
\toprule
Forecaster & Brier $\downarrow$ & $\Delta$ base & Log loss $\downarrow$ \\
\midrule
Untrained Qwen3-8B & .1287 & --- & .5003 \\
Trained, seed 42 & .1164 & -.0123 & .4058 \\
Trained, seed 43 & .1168 & -.0119 & .4069 \\
Trained, seed 44 & .1161 & -.0126 & .4146 \\
\midrule
Trained mean & .1164 & -.0123 [-.01948,-.00602] & .4091 \\
\bottomrule
\end{tabular}
}
\end{minipage}
\end{center}

\begin{center}
\begin{minipage}{.9\linewidth}
\centering
\captionof{table}{\textbf{Qwen3-4B and Llama held-out comparisons.}
The Qwen3-8B result is immediately above. $\Delta$ columns are trained
mean minus comparator, so negative values favor training. Parse is the minimum
coverage across base and trained endpoints; Llama's 12 unparsed base outputs
receive Brier loss one.}
\label{tab:exp3b-model-roster-results}
\scriptsize
\setlength{\tabcolsep}{3.5pt}
\resizebox{\linewidth}{!}{%
% generated by polymarket/scripts/render_model_extension_table.py -- do not edit
\begin{tabular}{lrrrrrrrrr}
\toprule
Model & Base & Seed 42 & Seed 43 & Seed 44 & Trained mean & $\Delta$ base [95\% CI] & $\Delta$ market [95\% CI] & $\Delta$ Platt [95\% CI] & Parse \\
\midrule
Qwen3-4B & .1316 & .1142 & .1140 & .1194 & .1159 & -.0157 [-.02528,-.00667] & .0022 [-.00264,.00721] & .0036 [-.00125,.00874] & 1.000 \\
Llama-3.1-8B & .1520 & .1151 & .1149 & .1151 & .1150 & -.0370 [-.05376,-.02258] & .0013 [-.00004,.00371] & .0027 [.00067,.00539] & .988 \\
\bottomrule
\end{tabular}

}
\end{minipage}
\end{center}

Qwen3-4B replicates the base-relative improvement: its trained mean is $.11589$
versus $.13160$ for base, a difference of $-.01572$
$[-.02528,-.00667]$. Llama-8B improves from $.15203$ to $.11502$, a difference
of $-.03701$ $[-.05376,-.02258]$; all three trained Llama endpoints parse.

The latest pre-cutoff market price obtains Brier $.11369$, and logistic
recalibration of its logit fit only on the training split obtains $.11232$.
Neither comparator informs language-model training. The trained-minus-market
interval spans zero for every model. Qwen3-4B is also statistically
indistinguishable from the calibrated price, but Qwen3-8B and Llama-8B have
higher error by $+.00409$ $[+.00098,+.00785]$ and $+.00269$
$[+.00067,+.00539]$, respectively. Thus the complete roster shows
base-relative improvement rather than superiority to market-based forecasts.

\subsection{Scale and Architecture Comparison}
\label{app:exp4-scale}

The scale comparison uses all 318 eligible 2026-window tasks whose identifiers and
event, series, and normalized question-template families are absent from the
training, development, and 1,024-market evaluation sets. The holdout contains 227
event families and 288 normalized templates; selection does not use outcomes or
market probabilities.

Qwen3 checkpoints at 4B, 8B, and 14B use three seeds, a common reward, a
4,800-presentation schedule, and the same development rule. Each untrained base
and trained adapter receives five non-thinking draws per task at temperature
$.7$, top-$p$ $.8$, and top-$k$ 20. All 19,080 Qwen draws satisfy the strict JSON
contract.

\begin{center}
\begin{minipage}{.94\linewidth}
\centering
\captionof{table}{\textbf{Exact values behind the main-text scale figure.}
Lower Brier is better. Trained means average seeds 42--44. Missing five-draw task
groups receive Brier loss one.}
\label{tab:exp4-qwen-scale-results}
\scriptsize
\setlength{\tabcolsep}{3.5pt}
\resizebox{\linewidth}{!}{% Generated by paper/generate_exp4_scale_figure.py; do not edit manually.
\begin{tabular}{lrrrr}
\toprule
Checkpoint & Base & Trained mean & Trained $-$ base & Trained parse \\
\midrule
Qwen3-1.7B & 0.2253 & 0.1248 & -0.1005 & 100.0\% \\
Qwen3-4B-Instruct-2507 & 0.1522 & 0.1366 & -0.0156 & 100.0\% \\
Qwen3-8B & 0.1369 & 0.1286 & -0.0083 & 100.0\% \\
Qwen3-14B & 0.1281 & 0.1270 & -0.0012 & 100.0\% \\
Llama-3.2-3B & 0.2189 & 0.1367 & -0.0822 & 100.0\% \\
Llama-3.1-8B & 0.1642 & 0.1264 & -0.0377 & 100.0\% \\
\bottomrule
\end{tabular}
}
\end{minipage}
\end{center}

The trained means are $.13664$, $.12859$, and $.12695$ at 4B, 8B, and 14B;
the corresponding bases are $.15224$, $.13686$, and $.12812$. The 14B-minus-8B
trained contrast is $-.00164$ $[-.00332,-.00024]$, and the 14B-minus-4B
contrast is $-.00969$ $[-.01869,-.00148]$. The 14B trained-minus-base contrast
is $-.00117$ $[-.00298,+.00016]$.

On this holdout, contemporaneous Polymarket prices have Brier $.12655$.
Trained 14B minus market is $+.00040$ $[-.00028,+.00139]$, so the interval
crosses zero. Thus 14B has lower trained error than 8B while its base-relative and
market-relative differences remain unresolved.

Llama-3.1-8B uses the same holdout and five-draw endpoint. All 4,770 trained draws
parse; nine of the 1,590 base draws do not, causing seven base task groups to
receive Brier loss one. Trained seed Briers are $.12638$, $.12629$, and
$.12662$, for a mean of $.12643$ versus $.16417$ for base. Trained minus base is
$-.03774$ $[-.05860,-.01946]$; trained minus Polymarket is
$-.00012$ $[-.00046,+.00012]$. The latter interval crosses zero.

Qwen3-1.7B and Llama-3.2-3B preserve the same training/development data, three
seeds, 4,800-presentation schedule, five stochastic draws, and holdout.
Qwen3-1.7B seed Briers are $.12413$, $.12610$, and
$.12415$ (mean $.12479$) versus $.22531$ for base; trained minus Polymarket is
$-.00175$ $[-.00644,+.00265]$. Llama-3.2-3B seed Briers are $.13197$, $.13666$,
and $.14149$ (mean $.13671$); trained minus Polymarket is $+.01016$
$[+.00121,+.02060]$. The pinned public BF16 Llama mirror is recorded because the gated upstream
repository was unavailable to the workspace account. Its unconstrained base
initially produced only 4/318 complete five-draw groups at the 128-token cap;
an output-cap-only rerun reached 21/318. A disclosed, holdout-informed syntax
recovery therefore repeated the full base endpoint with identical prompts and
sampling settings while constraining output to a JSON object containing one
numeric probability. All 1,590 recovered draws parse. Base Brier is $.21887$,
and trained minus base is $-.08216$ $[-.11062,-.05367]$. Because the syntax
constraint was introduced after observing the failures and differs from the
trained endpoints, this base-relative contrast is descriptive. Both trained
models also parse every draw.

\paragraph{Hosted frontier references.}
The API panel evaluates the same 318-task holdout. It is categorical rather than a parameter-scale
curve: provider parameter counts are unavailable or not comparable, and the
APIs do not expose identical decoding controls. DeepSeek-V4-Pro uses the common
temperature $.7$/top-$p$ $.8$ stochastic request. GPT-5.6-Sol uses provider-native
low reasoning without user-set temperature or top-$p$; Claude Opus~4.8 likewise
uses provider-native sampling. Each endpoint receives five independent calls
per task and uses the same strict parser, five-draw mean, Brier scoring, and
family-clustered 5,000-repetition interval procedure as the local panel.

DeepSeek parses all 1,590 draws and has Brier $.12660$, differing from the crowd by
$+.00005$ $[-.00105,+.00102]$. Claude parses all draws and has Brier $.12927$,
a difference of $+.00272$ $[-.00049,+.00614]$. GPT receives all 1,590
responses but six are malformed; two task groups therefore receive Brier loss
one, and task-level coverage is 99.37\%.
Its Brier is $.13335$, a difference of $+.00680$
$[+.00016,+.01635]$.

Kimi uses a 1,024-token output budget. Thirteen transport failures are retried
without replacing the 1,577 successful responses; all 1,590 final draws parse.
Its Brier is $.12682$, differing from the crowd by $+.00027$
$[-.00046,+.00109]$. Hosted results provide performance context, not evidence
for a common scale law or an identical-decoder comparison.

\subsection{Qwen3-4B Evidence-Updating Evaluation}
\label{app:exp3b-update}
This evaluation applies the Experiment~1 update prompt and all nine fixed packets
for each of its 100 markets (three pro-$H_1$, three
anti-$H_1$, and three orthogonal). The packet and initial-context SHA-256 digests
are \texttt{6741be0a\ldots8796d36} and
\texttt{c550cf05\ldots0c098}. Every condition receives the same canonical first
structured Qwen~2.5-7B forecast for each market, so this paired analysis isolates
conditional revision rather than web retrieval or initial event-model
construction. Search and thinking are disabled; decoding is greedy with a
1,500-token output cap. The longest prompt is 1,480 tokens against a 6,144-token
context limit. These contrasts characterize revision behavior and are separate
from the held-out forecasting evaluation.

The paired evaluation produces all 900 packet-conditioned revisions for the base
and each final adapter, with 100\% parse coverage and ICS equal to one. Base
EHC/HFC is $.909$ on 591 eligible directional updates. Seeds 42, 43, and 44 obtain
$.917$, $.922$, and $.918$ (592, 588, and 591 eligible updates), for a
three-seed mean of $.919$. The paired market-bootstrap trained-minus-base EHC
difference is $+.0101$ with 95\% interval $[+.0004,+.0207]$; HFC has the same
estimate and interval $[+.0009,+.0210]$. Training reduces mean absolute
directional revision by $.52$ percentage points $[-.70,-.36]$ and orthogonal
revision by $.22$ points $[-.38,-.06]$; the directional-minus-orthogonal gap
narrows by $.30$ points $[-.52,-.09]$. The sensitivity ratio changes from
$4.64\times$ to $4.83\times$ ($+.19$ $[-.07,+.48]$). These components show
improved directional consistency alongside smaller revisions overall; they do
not indicate an increase in absolute packet selectivity.
\begin{center}
\begin{minipage}{.82\linewidth}
\centering
\captionof{table}{\textbf{Paired evidence-updating contrasts.}
Movement rows are mean absolute probability revisions in percentage points;
intervals resample the 100 markets.}
\label{tab:exp3b-update-seed-results}
\scriptsize
\setlength{\tabcolsep}{4pt}
\resizebox{\linewidth}{!}{%
% Paired Experiment 1 reapplication contrasts, 2026-08-23.
% Source: reports/exp3b_exp1_reapplication_j30535058/movement_component_analysis.json.
\begin{tabular}{lrr}
\toprule
Measure & Trained $-$ base & 95\% interval \\
\midrule
Directional consistency (EHC) & $+.0101$ & $[+.0004,+.0207]$ \\
Headline-forecast consistency (HFC) & $+.0101$ & $[+.0009,+.0210]$ \\
Directional movement & $-.52$ pp & $[-.70,-.36]$ \\
Orthogonal movement & $-.22$ pp & $[-.38,-.06]$ \\
Directional $-$ orthogonal movement & $-.30$ pp & $[-.52,-.09]$ \\
Sensitivity ratio & $+.19\times$ & $[-.07,+.48]$ \\
\bottomrule
\end{tabular}

}
\end{minipage}
\end{center}

\section{Experiments 3--4: Shared Training Configuration and Runtime}
\label{app:exp3-complete-training-setup}

Tables~\ref{tab:exp3-oat-common}--\ref{tab:exp3-runtime} record the effective
OAT/Dr.~GRPO~\citep{liu2024oat} arguments for the reported Coin City and Polymarket studies, including framework defaults omitted from the command line.
Execution manifests hash the data, trainer, launcher, evaluator, prompt
renderer, and each cell's environment.

\paragraph{Step accounting.}
A ``step'' in the figures and earlier tables is one rollout-and-update round rather than one
AdamW call.  Each full run consumes 16 prompts per round and samples eight responses
per prompt (128 trajectories).  With one learner GPU, per-device microbatch one, one
PPO epoch, and gradient accumulation 16, OAT performs eight AdamW parameter updates per
round. Thus 4,800 prompt presentations give 300 reported rounds and 2,400 AdamW
updates. The command line supplied a .03 warmup ratio, but OAT selected
Transformers' \texttt{constant} scheduler, which does not use the warmup argument.
The effective learning rate was therefore $10^{-6}$ from the first optimizer
update. We use ``round'' below to avoid counting the eight within-round parameter
steps as independent data exposures.

\begin{center}
\begin{minipage}{\linewidth}
\centering
\captionof{table}{\textbf{Shared Dr.~GRPO implementation parameters.} These
values apply to the Coin City and Polymarket studies; the study-specific
tables below record exceptions.}
\label{tab:exp3-oat-common}
\scriptsize
\setlength{\tabcolsep}{3pt}
\begin{tabular}{p{.29\linewidth}p{.66\linewidth}}
\toprule
\textbf{Quantity} & \textbf{Effective value} \\
\midrule
Trainable parameters & Frozen base model; LoRA rank 32, $\alpha=64$, dropout 0,
no bias; targets \texttt{q\_proj}, \texttt{k\_proj}, \texttt{v\_proj},
\texttt{o\_proj}, \texttt{down\_proj}, \texttt{up\_proj}, and
\texttt{gate\_proj}; no quantization. \\
Algorithm & OAT Dr.~GRPO (\texttt{critic\_type=drgrpo}); eight-response group
advantage $r_i-\bar r$ with no group-standard-deviation division; policy loss is
token-summed with the generation limit as a fixed normalizer (no response-length
normalization). \\
Policy objective & PPO ratio clipping 0.2; truncated importance-sampling cap 2.0;
one PPO epoch; KL coefficient $\beta=0$; no critic/value model; reward whitening,
advantage whitening, and REINFORCE mode all off. \\
Optimizer & PyTorch AdamW \citep{loshchilov2019decoupled} (substituted for
DeepSpeed fused/CPU Adam); learning rate
$10^{-6}$; $\beta_1=.9$, $\beta_2=.95$; weight decay 0; maximum gradient norm 1.0. \\
Schedule & Constant $10^{-6}$ learning rate from the first optimizer update; the
supplied .03 warmup ratio is ignored by Transformers' \texttt{constant} scheduler;
one prompt epoch; 4,800 prompt presentations $=$ 300 rollout-and-update rounds
$=$ 2,400 AdamW steps. \\
Batching & 16 prompts per round; 8 sampled responses per prompt; 128 trajectories;
training microbatch 1; gradient accumulation 16; one learner GPU. \\
Training sampling & Temperature 1.3; top-$p=1$; top-$k=-1$; one independently
sampled completion at a time, repeated eight times per prompt. \\
Truncation and stopping & Model EOS only (no external stop string); a completion
ending at the length limit receives reward zero.  Missing or invalid required JSON
also receives reward zero. \\
Precision and memory & bfloat16; FlashAttention~2; gradient checkpointing;
DeepSpeed ZeRO-2 \citep{rajbhandari2020zero}; reference-model offload; vLLM prefix
caching \citep{kwon2023vllm}; dropout disabled. \\
Seeds and shuffling & Training seeds 42, 43, and 44.  Each dataset is deterministically
shuffled once, then OAT shuffles prompt batches independently under the training seed. \\
Checkpointing and logging & Evaluation includes step zero and the final round;
intermediate evaluation at the setting-specific frequency; one final LoRA adapter retained; metrics
logged each round; Weights \& Biases disabled. \\
\bottomrule
\end{tabular}
\end{minipage}
\end{center}


\subsection{Coin City A-to-B Factorial}
\label{app:exp3-coin-structural-hparams}

Table~\ref{tab:exp3-coin-structural-hparams} records the shared configuration
for Experiment~3. The main analysis uses five-draw stochastic
round-300 evaluations from the three Qwen3-8B episode-matched and three
population-prior seeds. The appendix reports the Qwen3-4B stochastic
replication, its three-seed structureless diagnostic, and the Llama-8B
architecture comparator. The
table supplements the shared optimizer details in Table~\ref{tab:exp3-oat-common}; where the two tables differ, this
setting-specific table governs. The execution manifest stores the environment;
the protocol archive hashes the preflight manifest, data files,
generator, prompt renderer, scorer, and trainer entry point.

\begin{center}
\begin{minipage}{\linewidth}
\centering
\captionof{table}{\textbf{Evaluated configuration for Coin City
hidden-structure transfer.}}
\label{tab:exp3-coin-structural-hparams}
\scriptsize
\setlength{\tabcolsep}{3pt}
\begin{tabular}{p{.29\linewidth}p{.66\linewidth}}
\toprule
\textbf{Quantity} & \textbf{Effective value} \\
\midrule
Model roster & \texttt{Qwen/Qwen3-4B-Instruct-2507},
\texttt{Qwen/Qwen3-8B}, and \texttt{meta-llama/Llama-3.1-8B-Instruct}. \\
Chat formatting & Qwen3-4B uses the explicit ChatML wrapper; Qwen3-8B and
Llama-3.1-8B use their tokenizer templates, with Qwen thinking disabled. All
include an assistant-generation prefix. \\
Evaluated grid & 3 models $\times$ 2 reward arms
(episode-matched reward, stored under the legacy artifact key \texttt{causal}, and \texttt{population\_prior}) $\times$ 3 seeds
($42,43,44$) $=18$ full training runs. \\
Diagnostics & Qwen3-4B \texttt{structureless} at the same three seeds, plus one
untrained endpoint for each model. These six diagnostics are excluded from the primary
episode-matched-versus-prior estimand. \\
Training data & Coin City and Structure~A only; 4,800 prompts, balanced across
$k\in\{0,2,4,8\}$ and weak/strong target match. Two reference trajectories of
length 14 and one target trajectory of maximum length 8. \\
Comparator target construction & Episode-matched reward uses the episode truth.
Population-prior reward targets are a fixed no-self-match derangement within
each $k$; its target-vector multiset is exactly identical to episode-matched training.
Prompt strings are byte-identical across these two arms. \\
Held-out data & $2$ domains $\times2$ target structures $\times3$ hint arms
$\times4$ evidence depths $\times30$ episode seeds $=1,440$ rows. The complete
held-out dataset is byte-identical across all training arms. \\
Mechanism ranges & Evaluation $g\in[.22,.92]$; Structure~B
$\rho\in[.48,.74]$, $\phi\in[.18,.42]$, and $\lambda\in[.68,.94]$.
Training reference gains are $[.18,.46]$ and $[.68,.98]$. \\
Domain ranges & Coin City: baseline 50, support $[0,100]$, noise 1.6, driver
pool $\{-12,-8,-4,0,4,8,12\}$. Coin Harbor: baseline 180, support $[40,320]$,
noise 4.5, driver pool $\{-24,-16,-8,0,8,16,24\}$. \\
Interventions & Five shocks at horizons one and three; Coin City
$\{-12,-6,0,6,12\}$, Coin Harbor $\{-24,-12,0,12,24\}$. Ten joint forecasts
produce eight nonzero-minus-zero response contrasts. \\
Reward & $.4\,r_{\rm level}+.6\,r_{\rm response}$. Each component averages
$\max(0,1-|e|/S)$; level tolerance is 10\% of the domain's observable range
(10 Coin City points; 28 Coin Harbor units), and response tolerance is half one
shock unit (3 and 6, respectively). Invalid or length-truncated output receives
zero online reward. \\
Output contract & Exactly one JSON object with sole key
\texttt{forecasts} and exactly ten finite JSON numbers. The same strict parser
is used online and offline; no prose, extra/duplicate key, Boolean, non-finite
value, or wrong-length list is accepted. \\
Structured decoding & Syntax-only XGrammar GBNF in training, online evaluation,
endpoint evaluation, and base evaluation. It fixes object shape and arity only;
numeric values and their relationship to targets remain unconstrained. \\
Optimization & Frozen base plus LoRA rank 32, $\alpha=64$; OAT Dr.~GRPO;
PyTorch AdamW learning rate $10^{-6}$ with an effective constant schedule and no
warmup; $\beta=0$; one PPO epoch; gradient norm 1; ZeRO-2 and reference offload. \\
Batch and update accounting & 16 prompts per round; eight rollouts per prompt;
128 trajectories; microbatch 1; gradient accumulation 16; 4,800 presentations
$=$ 300 rollout-and-update rounds $=$ 2,400 AdamW steps. \\
Training generation & Temperature 1.3, top-$p=1$, top-$k=-1$, model EOS only;
192 new-token limit; prompt limit 2,304; total model context 3,072; vLLM prefix
caching enabled. \\
Precision / memory & BF16, FlashAttention~2, gradient checkpointing, dropout
off, no quantization; vLLM GPU fraction .38; expandable-segments CUDA allocator. \\
Online evaluation & All 1,440 rows at temperature 0, one draw, top-$p=1$,
192-token limit, batch 120, at step zero and every 50 rounds including round 300. \\
Reported endpoint evaluation & Five draws per prompt at temperature .7 and
top-$p=.95$; trained decode seed is training seed $+20{,}260{,}818$ and the
base seed is 20,260,819. Temperature-zero checkpoints are retained only as
online monitoring artifacts and are not reported as endpoints. \\
Offline scoring & Level MAE, response MAE, response error normalized by the
blind 50/50 pathway mixture, response correlation, strict parse coverage, and
nearest behavioral pathway. A parse failure receives the full domain range as
both level and response error. \\
Variance and estimands & Five-draw outputs first average within episode.
Hierarchical intervals resample three training seeds, then paired episodes.
Primary: prior-minus-episode-matched response MAE for correct-hint Structure~B in Coin
Harbor, equal over $k$. Domain, structure, joint, hint, and hint-by-$k$ contrasts
form the factorial analysis. \\
Preflight validation & A 10-round run for every model and evaluated arm must
save an adapter, produce the full temperature-zero monitoring and five-draw
stochastic evaluation sets, achieve
at least .95 strict parse rate under each decoder, and exhibit nonconstant finite
reward. The 17-check data preflight must also pass before estimating the factorial. \\
Compute & Training and trained-checkpoint evaluation use two NVIDIA A6000 48-GB
GPUs (one actor, one learner), 8 CPU cores, and 56--100 GB allocated host memory.
Untrained base evaluation uses one A6000, 8 CPU cores, and 24 GB host memory.
All runs load models from the same pinned local cache. \\
\bottomrule
\end{tabular}
\end{minipage}
\end{center}

\begin{center}
\begin{minipage}{\linewidth}
\centering
\captionof{table}{\textbf{Runtime environment for the reported training runs.}
Package versions come from the job execution environment.}
\label{tab:exp3-runtime}
\scriptsize
\setlength{\tabcolsep}{3pt}
\begin{tabular}{p{.29\linewidth}p{.66\linewidth}}
\toprule
\textbf{Component} & \textbf{Version / configuration} \\
\midrule
Python / CUDA build & Python 3.10.20; PyTorch 2.6.0 built for CUDA 12.4. \\
Training stack & OAT 0.2.4 at Git revision
\texttt{869706620a35ec5304fa80a0\allowbreak{}c74f0c566cc2bc57}; DeepSpeed 0.16.8; PEFT 0.15.2. \\
Generation stack & Transformers 4.51.3; vLLM 0.8.4; FlashAttention 2.7.4.post1. \\
Data stack & Datasets 2.16.1; Accelerate 1.14.0. \\
Allocator / logging & \texttt{PYTORCH\_CUDA\_ALLOC\_CONF=expandable\_segments:True};
Weights \& Biases disabled; stdout/stderr and per-round metrics retained with each
run. \\
\bottomrule
\end{tabular}
\end{minipage}
\end{center}


\section{Experiments 3--4: Reproducibility and Artifact Map}
\label{app:exp3-repro}

The artifact roots \path{exp3_training_transfer/coin_city_structural/}
and \path{exp3_training_transfer/polymarket/scripts/} contain the generators, prompt builders,
trainers, evaluators, and analysis code for the Coin City and historical-market
studies. Execution ledgers record effective hyperparameters, Git revisions,
adapter identities, and SHA-256 hashes of data and executable protocol files;
result renderers verify these fields before aggregation. For Coin City,
\path{make_qwen3_8b_endpoint_outputs.py} validates the seven five-draw files
(50,400 scored draws) that support Figure~\ref{fig:exp3-transfer-results}.
\path{make_paper_outputs.py} validates the complete 24-endpoint roster and all
48 recorded endpoint files before writing the machine-readable result and the
stochastic appendix comparator table. Temperature-zero files are checked as
execution provenance but are excluded from manuscript estimates. For the
historical study,
\path{analyze_exp1_reapplication.py} validates four matched 900-record files and
computes the paired-market direction, sensitivity, and movement-component
contrasts. The Polymarket manifests preserve sample membership, prompt and
adapter hashes, parse coverage, and bootstrap settings.
